El Quetzal-2 es un proyecto académico desarrollado por estudiantes y personal de la Universidad del Valle de Guatemala la cual pone a prueba una computadora a bordo (OBC por sus siglas en inglés) diseñada en UVG capaz de ejecutar un modelo de inteligencia artificial para identificar nubes en imágenes satelitales \cite{quetzal_2}. Este proyecto es importante no solo porque representa un avance tecnológico en Guatemala, sino también por su impacto dentro de la juventud del país, inspirando a futuros científicos, ingenieros e innovadores \cite{quetzal_1}.

La misión espacial es un CubeSat 2U, 10x20x10cm y con un peso máximo de 4kg \cite{cubesat_2022} y cuenta con 4 objetivos específicos que busca completar, llamadas cargas útiles (\textit{Payloads} en inglés). Para este proyecto nos interesa específicamente el \textit{inflight software} de la OBC secundaria de Quetzal-2 y su interacción con \textit{Payload MILO}, el subsistema encargado del reconocimiento de nubosidad y la toma de fotografías en el espacio \cite{quetzal_2}. Es importante notar que estos 4 objetivos de misión son independientes de los 3 objetivos específicos de este trabajo profesional detallados en la sección de objetivos.

Las OBCs cumplen un papel clave dentro de los subsistemas del satélite puesto que son las encargadas de manejar todas las tareas, intercambio de información entre módulos y la colecta de información sobre los demás subsistemas (\textit{housekeeping}) antes de la conexión con la estación en tierra \cite{lwabanji_wilkinson_biermann_bellville_2013}.

Algunos de los subsistemas que debe manejar una OBC son:
  - Un sistema de poder (EPS por sus siglas en inglés).
  - Un sistema que determina la altidud (ADCS por sus siglas en inglés).
  - Un sistema de comunicación utilizando radiofrecuencias.
  - Una o más cargas útiles, por ejemplo una cámara o transmisor de señales junto con su controlador.

La inicialización, intercomunicación y gestión de todos estos subsistemas recae sobre la OBC \cite{akhtar_2012}. Debido a las altas limitaciones energéticas en satélites de esta escala se debe tener especial cuidado con el presupuesto de energía que se le puede brindar a cada uno de los subsistemas. De esta necesidad nacen los modos de operación, los cuáles restringen las funcionalidades del satélite según la tarea para la que se quiere que se especialice en ese momento.

Este trabajo profesional busca desarrollar un sistema de \textit{software} para la OBC diseñada localmente (\textit{in-house}) capaz de operar la carga útil MILO, gestionar modos de operación y ejecutar un mecanismo de \textit{handover} con la computadora principal, validado en un entorno controlado terrestre. Para el desarrollo de este sistema se utiliza una metodología propia del laboratorio espacial UVG, basada en la metodología de ingeniería de sistemas de la NASA \cite{nasa_systems_engineering}.
